\documentclass[11pt]{article}

% Use Helvetica font.
%\usepackage{helvet}

%\usepackage[T1]{fontenc}
%\usepackage[sc]{mathpazo}
\usepackage{color}
\usepackage{amsmath,amsthm,amssymb,multirow,paralist}
%\renewcommand{\familydefault}{\sfdefault}

% Use 1/2-inch margins.
\usepackage[margin=0.8in]{geometry}
\usepackage{hyperref}

\begin{document}

\begin{center}
{\Large \textbf{Formula 1 Race Algorithm Prediction}

\vspace{10pt}

Project Midterm Report}

\vspace{10pt}

\textit{Team Members: Sandeeptha Madan, Evan Sivets}
\end{center}

% \linethickness{1mm}\line(1,0){498}

%%%%%%%%%%%%%%%%%%%%%%%%%%%%%%%%%%%%%%%%%%%%%%%%%%%%%%%%%%%%%%%%%%%%%%%%%%%%%%%

%%%%%%%%%%%%%%%%%%%%%%%%%%%%%%%%%%%%%%%%%%%%%%%%%%%%%%%%%%%%%%%%%%%%%%%%%%%%%%%

\begin{abstract}
Our task in this project is to predict the winners of the races in the upcoming Formula 1 season. Our motivation is a strong interest in the sport and in seeing our work come to life, to see how accurate our results are. We have done projects related to Formula 1 in the past and want to build on that to explore something new. We will use data from previous years on all racers, tracks, and teams, and feed it into the models we used in class to predict this year's race outcomes accurately. So far, we have found that the SVM Linear model has been the best. We also did a practice test based on a previous year, and it correctly predicted the winner, Leclerc, at the 2022 Bahrain GP.
\end{abstract}

\section{Introduction (Done/Look Over)}

Formula 1 is the highest level of international auto racing, featuring 24 racing weekends across international circuits. This means there will be a different track to race on each event weekend. Different teams are racing as well, each with two racers. These teams each have different vesicles, all very similar to one another, making some vehicles better than others. This project will use past data from race tracks, teams, racers, and finishes to predict the podium placers for every race this coming year. 

The datasheet that we will be using is "Formula 1 World Championship - Kaggle." We will be using data from 2014-2024, as pre-2014 regulations differ in ways that make the constructor/driver strength non-comparable to previous eras. The models we will be using to test these data sets are Logistic Regression (L2), Logistic Regression (L1), Linear SVM, RBF SVM, and Decision Tree. With these models, the following tasks will be to take into account the starting grid position, qualifying classification position, drivers' podium rate across all prior races this season, constructors' podium rate across all prior races, Drivers' average grid position across all prior races, and drivers' podium rate minus their teammates' podium rate.

The expected outcomes of this project are that the SVM Linear will be the best model to use, as linear models should follow the data, and that SVM will be the best margin classifier to group everything. We do not expect to predict every race correctly, but we do expect the results to be close to what was calculated based on statistics the drivers and teams have shown from past races on the tracks that they will see this season.

\section{Related Work (Done 1, could add a second in the future)}

There has been a similar study to ours by Spencer Staub at The George Washington University in Autumn 2022, titled "Formula One Race Prediction Model." He took a look at all the features in our data, such as drivers' previous races at tracks, teams' past results, qualifying win rates, vehicle engineering, and drivers' past finishes. They used similar models we used in Logistic Regression and SVM, plus one we did not use: XGBoost. Spencer also took a look at the improvement results from the previous season, which had an accuracy of 77 percent, suggesting the stats would be in our favor and something to take into account when comparing our work with past stats. In their future work, they plan to use lap data for testing, which we currently have in our project, to help with predicting the outcomes of future races.

\section{Your methods name (To Do)}

In this section, you will introduce your machine learning
algorithms used in this part.

If you are using existing methods on a \textbf{new dataset}. You
can talk about ``dataset collection'', ``data pre-processing'',
``Feature selection'', ``Choices of hyper-parameter'',
``Evaluation metrics'', etc.

If you propose some \textbf{new methods}. You can talk about
``Drawbacks of existing methods'', ``Your proposed methods'',
``Your models based on your methods'', etc.

These are possible subsections in this part. You don't need to
include all of them. I want to see the contents that you spend
time on and have some insights.

This part is one of the most important sections. The length of this
section should be about 1.5 pages.

\section{Preliminary Results(Done/Look over)}

To help show what we are wanting to do with this project, we started off by predicting a race that had already happened to see what the results would be. We feed the model features that we already know ahead of time, grid position, qualifying result and how the driver and constructor have been performing that season. The race we decided to use was 2022 Bahrain GP where Leclerc won. The model used the following data and predicted him to win, which he did. 

We used our six models stated in the introduction to run tests on two Sets, Validation(2022) and Test(2023-2024). With these sets we found the accuracy, precision, recall, F1 and ROC-AUC of each model. With our Validation set SVM linear had the highest percentage in each of the 5 categories tested, while Decision tree being the worst in about every category. With the Test set SVM linear had the highest ROC-AUC while other models had the highest percentages else where. With these high percentages this means that the models can catch almost every podium finish accurately. 

In one of our other tests we took two models, Decision Tree and Logistic Regression(L2) -|Coefficients|, and tested what is the importance of each feature. Decision Tree said the that qualifying postion as the most importance as an 80 percent importance value, with teammate podium rate difference being the least important. L2 stated that the grid was the most important, and matching the Decision tree in having the same feature with the least importance. Although these models had different importance values across features, they each had the same top three and bottom three just in different order.

In another test we made a heatmap on how the features correlate with one another. This map will help us determine the most important features. These would be grid position, qualifying position and, driver/constructor podium rates. It also helps to take into account what features can be avoided and are redundant, for example teammate podium rate diff.



\section{Future plan(Done/Look Over)}

A key next step for this project is to evolve it from a solid classification model into a more rich and realistic race-prediction system. While the current model performs well using only pre-race features, future efforts should focus on enhancing generalization and practical utility. One way to achieve this is by implementing season-by-season validation to ensure the model remains reliable across different years, particularly given the frequent shifts in team performance in Formula 1. 

Expanding the feature set to include recent race form, circuit-specific performance, and reliability indicators would provide a more comprehensive representation of race conditions. Additionally, since the model currently prioritizes high recall at the expense of precision, refining probability thresholds or adopting a ranking-based approach—where drivers are ordered by the likelihood of achieving a podium finish—would make the predictions more meaningful.

Incorporating advanced models, such as ensemble methods like Random Forests or XGBoost, could further enhance performance. Finally, presenting the predictions in a more user-friendly format, such as a dashboard or race preview tool, would increase the project’s impact and make it more appealing as both a technical and applied machine learning solution.

\section{References (To Do)}

https://datasci.columbian.gwu.edu/sites/g/files/zaxdzs4746/files/2023-02/f1-final-presentation.pdf

https://www.kaggle.com/datasets/rohanrao/formula-1-world-championship-1950-2020

\end{document}
